\section{Method: Bottleneck-Aware Diagnostic Adaptation}\label{sec:method}

\subsection{Overview}

The \framework{} has three moving parts that compose into a single
inference-time pipeline. (1)~A \emph{calibration trace} is run once per
benchmark slice: a small fixed set of candidate routes is evaluated on rows
$0$--$199$ of each dataset, producing per-method accuracy, gold-passage
retrieval rate, gold-conditional accuracy, per-row disagreement across
routes, parse/error counts, and average LLM-call cost. (2)~A
\emph{bottleneck label} is assigned to the slice from a fixed signal set
(Section~\ref{sec:method:signals}); each label corresponds to a known
failure mode and a known matching intervention. (3)~At inference, the
\controller{} performs \emph{per-row routing} to the cheapest intervention
in our portfolio whose preconditions match the slice's bottleneck label
and the row's diagnostic signature. The full pipeline is therefore
\emph{calibrate $\rightarrow$ label $\rightarrow$ route}, with retrieval,
reranking, and embedding shared across all routes. We emphasize at the
outset that the \controller{} is a rule and evidence-summary policy over
the calibration trace, \emph{not} a learned router: there is no policy
network, no reflection-token training, and no question-complexity
regression. The mapping from signals to intervention is auditable in a
small table and is reported in Section~\ref{sec:method:controller}.

\subsection{Intervention portfolio: the method ladder}\label{sec:method:ladder}

The candidate routes form an inherited \emph{method ladder} in which each
step adds exactly one capability over the step below it, so that a
negative result at a higher rung points at a specific intervention rather
than at an under-specified bundle. Figure~\ref{fig:method_ladder} shows
the ladder; Table~\ref{tab:method_costs} lists the LLM-call cost of each
route. The diagnostic \controller{} chooses among these per row; the
ladder is not run end-to-end on any single question.

\begin{figure}[t]
\centering
\includegraphics[width=\columnwidth]{16_method_ladder_flowchart.png}
\caption{Method ladder. Each step adds one capability so negative results
remain interpretable. The diagnostic controller picks per row.}
\label{fig:method_ladder}
\end{figure}

\begin{table*}[t]
\centering
\caption{Method ladder used as the intervention portfolio. Counts refer
to LLM generation calls per question; retrieval, embedding, and reranking
are not counted.}
\label{tab:method_costs}
\small
\begin{tabularx}{\textwidth}{lcY}
\toprule
Route & Calls & Capability added \\
\midrule
\method{rag\_simple} & 1 & baseline dense+rerank retrieval \\
\method{rag\_state\_filter} & 1 & jurisdiction-scoped metadata filter \\
\method{snap\_only\_in\_final} & 2 & snap reasoning anchors final synthesis \\
\method{rag\_rewrite} & 2 & legal query rewrite before retrieval \\
\method{rag\_hyde} & 2 & HyDE hypothetical passage as retrieval handle \\
\method{rag\_snap\_hyde\_2call} & 2 & fixed \hyre{} (snap-conditioned passage) \\
\method{adaptive\_snap\_hyre\_v2} & 2 & BarExam route: rewrite/HyRE selector \\
\method{adaptive\_snap\_hyre\_housing\_verifier} & 2 & state filter + yes/no verifier \\
\method{adaptive\_snap\_hyre\_diverse} & 2 & diverse \hyre{} passages for option spread \\
\method{adaptive\_snap\_hyre\_disagreement\_majority\_prior} & 2 & disagreement arbitration with majority prior \\
\method{adaptive\_snap\_hyre\_frontier} & 2 & frontier replay across HyRE-family rows \\
\bottomrule
\end{tabularx}
\end{table*}

Every route in Table~\ref{tab:method_costs} draws from the same dense
encoder, BM25 hybrid pool, cross-encoder reranker, and prompt sanitizer;
they differ only in (a)~what text is used as the retrieval handle, (b)~the
metadata predicate, if any, and (c)~the final-synthesis prompt. This
isolation is what makes the ladder diagnostic: a route winning over its
immediate neighbor identifies the single capability that closed the gap.

\subsection{Snap-HyRE: snap-conditioned hypothetical reasoning}\label{sec:method:snaphyre}

\hyre{} is the central new route in the portfolio and the route most
sensitive to the active bottleneck. In a single LLM call, the model is
prompted to produce two clearly labeled blocks: a \texttt{\#\#~Reasoning}
\emph{snap answer}, in which it commits to a candidate answer based on
parametric knowledge alone; and a \texttt{\#\#~Passage} block stating the
controlling rule, doctrine, statutory test, fact pattern, or legal
principle relevant to resolving the question. The \texttt{\#\#~Passage}
block is then embedded with the same dense encoder used to embed the
corpus and used as the retrieval handle for dense+BM25 hybrid retrieval
and cross-encoder reranking. The final synthesis call sees the retrieved
corpus passages plus the original question; the snap answer itself is
\emph{not} handed to final synthesis verbatim, only mediated through the
passages that the \texttt{\#\#~Passage} block successfully retrieved.

Two distinctions matter. Against \emph{standard HyDE}
\citep{gao2023hyde}, the hypothetical document is conditioned only on the
question and is asked to look like \emph{an answer}; against \hyre{}, the
hypothetical document is conditioned on a snap-committed answer and is
asked to state \emph{the rule that resolves the question}. The
architectural justification is brief: an answer-anchored prompt forces
the passage toward rule-like, retrieval-shaped text (``The rule against
perpetuities voids any interest that may not vest within twenty-one years
of a life in being'') rather than toward a paraphrase of the question
stem, which is the failure mode of pure HyDE on multiple-choice legal
items where the stem already carries most of the lexical surface.

We do \emph{not} claim \hyre{} is universally better than baseline RAG.
The calibration evidence in Section~\ref{sec:analysis} of the paper shows fixed \hyre{} underperforming
baseline on HousingQA ($51.5\%$ vs.\ $60.5\%$) and CaseHOLD ($72.0\%$
vs.\ $73.0\%$); the corresponding HyRE-only controls are weaker still
($50.0\%$ on HousingQA, $71.5\%$ on CaseHOLD). \hyre{} is therefore
positioned as a \emph{route} inside the portfolio: it lifts BarExam,
where the active bottleneck is query/legal-reasoning formulation, but
loses on slices where the active bottleneck is jurisdiction scope or
answer-option conversion. This is the precise sense in which the
controller's gain is bottleneck-explained: \hyre{} is selected where it
helps and withheld where the negative controls predict it will hurt.

\subsection{Diagnostic signals and bottleneck labels}\label{sec:method:signals}

For each candidate route on a calibration slice we record a compact
signal vector: route accuracy and the paired delta over the matched
baseline; gold-passage retrieval rate and gold-conditional accuracy,
where the dataset exposes gold passages; per-row disagreement (the
fraction of rows on which two routes give different predicted answers,
and the row-level oracle ceiling of picking the best of the two);
unparseable-answer and engine-error counts; and average LLM calls per
question. The same signals are computed for a top-$1$ vs.\ top-$5$
retrieval-depth ablation on the baseline, since the depth signature is
the cheapest discriminator between a depth bottleneck and a parametric
floor.

The eight bottleneck labels in Table~\ref{tab:bottleneck_taxonomy}
exhaust the failure modes we observed across the four legal benchmarks.
Each row pairs a diagnostic signal with the single matching intervention
that closes the gap, and names the benchmark on which that bottleneck
dominates. The labels are not mutually exclusive at the row level (a single
HousingQA row can be both jurisdiction-scoped and a false-positive yes), but
they are mutually exclusive as \emph{slice-level} labels by construction:
the controller picks the dominant signal.

\begin{table*}[t]
\centering
\caption{Bottleneck taxonomy. Each row pairs a diagnostic signal computed
from the calibration trace with the matching intervention from the method
ladder, and names the benchmark on which that bottleneck dominates.}
\label{tab:bottleneck_taxonomy}
\footnotesize
\begin{tabularx}{\textwidth}{p{0.15\textwidth}YYp{0.13\textwidth}}
\toprule
Bottleneck & Diagnostic signal & Intervention & Dominates on \\
\midrule
Retrieval depth & Catastrophic top-1 vs.\ top-5 drop on baseline & Larger $k$; smarter passage selection & Depth-sensitive datasets \\
Candidate-set saturation & Top-1 collapses; top-5/top-10 tie; gold-hit rises but answer flat & Restrict to small reranked candidate set & LegalBench-SCALR \\
Query / legal-reasoning formulation & Rewrite vs.\ HyRE asymmetry; snap reasoning lifts more than HyDE retrieval alone & \method{rag\_rewrite} or \method{adaptive\_snap\_hyre\_v2}; \emph{route} between them & BarExam \\
Metadata / jurisdiction scope & Empty-retrieval or near-zero state-match rate; explicit state filter beats generic top-$k$ & \method{rag\_state\_filter} & HousingQA \\
Statutory entailment / false-positive yes & Verifier route lifts over plain state filter; fixed \hyre{} underperforms baseline & Conservative yes/no verifier on top of state filter & HousingQA \\
Answer-option conversion & Gold-retrieval gains do not convert to answer gains (CaseHOLD gold-hit $16.0\%\!\rightarrow\!47.0\%$, EM $69.5\%\!\rightarrow\!72.0\%$, $p\!=\!0.44$) & Diverse \hyre{}; direct option-table is a clean negative & CaseHOLD \\
Method disagreement / candidate exposure & High disagreement across HyRE-family rows; row-level oracle headroom unrealized by simple ensembles & Disagreement arbitration or frontier replay selector & LegalBench-SCALR \\
Parametric floor / depth-flat & Top-1 $\approx$ top-5 $\approx$ \method{llm\_only}; snap-only matches RAG & Trust priors; use snap reasoning to anchor, not to retrieve & BarExam \\
\bottomrule
\end{tabularx}
\end{table*}

\subsection{The bottleneck-aware controller}\label{sec:method:controller}

The \controller{} is a hand-specified rule policy parameterized by the
calibration trace. Algorithm~\ref{alg:controller} gives the procedure end
to end.

\begin{algorithm}[H]
\caption{Bottleneck-aware diagnostic adaptation.}
\label{alg:controller}
\begin{algorithmic}[1]
\REQUIRE benchmark slice $S$ (rows $0$--$199$); intervention portfolio
$\mathcal{R}$ from Table~\ref{tab:method_costs}; bottleneck taxonomy
$\mathcal{B}$ from Table~\ref{tab:bottleneck_taxonomy} (label $\rightarrow$
matching-route subset).
\STATE \textbf{Calibration.} For each route $r \in \mathcal{R}$, run $r$
on $S$ and record the signal vector $s_S = (\text{acc}_r,
\text{gold-hit}_r, \text{cond-acc}_r, \text{pair-disagree}_{r,r'},
\text{parse-err}_r, \text{empty-retr}_r, \text{calls}_r)$.
\STATE \textbf{Label.} For each bottleneck $b \in \mathcal{B}$, compute the
signal magnitude $\sigma_b(s_S)$ defined in Appendix~\ref{app:taxonomy_examples}
(e.g., $\sigma_{\text{retrieval-depth}} = \text{acc}_{\text{top-5}} -
\text{acc}_{\text{top-1}}$, $\sigma_{\text{option-conv}} = \text{gold-hit}
- \text{cond-acc}$). Assign $\ell^\star = \arg\max_b \sigma_b(s_S)$.
\STATE \textbf{Route selection.} Among routes whose preconditions match
$\ell^\star$, pick the cheapest by average LLM calls:
$r^\star = \arg\min_{r \in \mathcal{R}_{\ell^\star}} \text{calls}_r$.
\STATE \textbf{Per-row downgrade.} For each held-out row $i$, if the snap
fails to parse, the state filter returns zero candidates, or two cheap
routes disagree on the answer, swap $r^\star$ for a slice-specific
fallback (typically \method{rag\_simple} or \method{rag\_state\_filter}).
\STATE \textbf{Inference.} Apply $r^\star$ (or the fallback) to each row
in the held-out slice; report accuracy and average calls.
\end{algorithmic}
\end{algorithm}

Three properties matter. (i)~The routing rule is a table lookup composed
with two $\arg$-ops; there is no policy network, no reflection-token
training, and no gradient flowing into the rule. (ii)~The slice-level
label $\ell^\star$ is fixed once calibration is done; only the per-row
downgrade step (step 4) consults the held-out row itself, and it only
downgrades to the fallback. (iii)~Every routing decision is auditable: a
reader can read $\sigma_b(s_S)$ off the calibration trace and reconstruct
the chosen route. This sits in deliberate contrast with Adaptive-RAG
\citep{jeong2024adaptive}, which learns a classifier over question
complexity, and Self-RAG \citep{asai2024selfrag}, which trains reflection
tokens to gate retrieval. Both prior controllers fix the
\emph{retrieval pipeline} and switch on a property of the question; our
\controller{} also switches on a property of the \emph{benchmark's
failure mode}, which is what enables it to withhold \hyre{} on HousingQA
and CaseHOLD where fixed \hyre{} is a negative control.

Empirically, this controller lifts macro accuracy from $71.9\%$ to
$77.9\%$ on the calibration slice at $1.30$ average LLM calls per
question, and from $71.5\%$ to $77.5\%$ on the held-out slice at $1.54$
calls, both at the sub-$2$-call budget set by the method ladder; the
per-benchmark decomposition and the cross-model sanity transfer are
reported in Section~\ref{sec:analysis}.
